\documentclass[11pt, ukrainian]{article}
\setlength\textwidth{190mm} \setlength\textheight{237mm}\setlength\voffset{-18mm} \setlength\hoffset{-33mm}\usepackage[T2A]{fontenc}
\usepackage[utf8]{inputenc}
\usepackage[ukrainian]{babel}
\begin{document}
{\Large \center  Варіанти завдань до лабораторної роботи №3\par 2021/22 н.р.\par \bigskip\bigskip}

«Зростаюча серія» – це послідовність чисел, що йдуть поспіль і неспадають, наприклад 2 2 2 5 7 7.

«Власна зростаюча серія» – це зростаюча серія, що зростає, наприклад 2  5  7.

«Спадна серія» – це послідовність чисел, що йдуть поспіль і незростають, наприклад 7 7 5 3 3 2.

«Власна спадна серія» – це спадна серія, що спадає, наприклад~7~5~3.

«Серія» – це зростаюча або спадна серія.

«Власна серія» – це власна зростаюча або власна спадна серія.

«Пилка» – це послідовність чисел, що йдуть поспіль і чергуються за неспаданням та незростанням, наприклад 2 2 2 5 3 7 4 6 5 9 (завершуватися може і неспаданням, і незростанням).

«Парна пилка» – це пилка, що завершується незростанням, наприклад 2 2 2 5 3.

«Непарна пилка» – це пилка, що завершується неспаданням, наприклад 2 2 2 5 3 7.

«Власна пилка» – це послідовність чисел, що йдуть поспіль і чергуються за зростанням та спаданням, наприклад 2 5 3 7 4 6 5 9 (завершуватися може і зростанням, і спаданням).

«Парна власна пилка» – це власна пилка, що завершується спаданням, наприклад 2 5 3 7 6.

«Непарна власна пилка» – це власна пилка, що завершується зростанням, наприклад 2 5 3 7.

\textbf{Увага!} 1. У пилці після неспадання обов’язково має бути незростання, а після незростання  – неспадання. Тобто послідовності 2 5 7 та 2 7 5 3 не утворюють пилку. Аналогічне стосується й власної пилки. 2. Пилка має починатися неспаданням, власна пилка – зростанням.

«Пік» – це послідовність чисел, що йдуть поспіль і спочатку зростають, а потім спадають, наприклад 2 3 5 4 3. 

«Воронка» – це послідовність чисел, що йдуть поспіль і спочатку спадають, а потім зростають, наприклад 5 3 2 3 4.

«Плато» – це послідовність чисел, що йдуть поспіль і спочатку зростають, потім повторюються, а потім спадають, наприклад 2 3 5 5 4 3.

«Яма» – це послідовність чисел, що йдуть поспіль і спочатку спадають, потім повторюються, а потім зростають, наприклад 5 3 2 2 3 4.

«Гора» – це послідовність чисел, що йдуть поспіль і спочатку не спадають, а потім не зростають, наприклад 2 2 3 3 5 7 7 6 6 4 3 3. «Гора» починається неспаданням, а закінчується незростанням. 

«Власна передгора» – це гора, що починається зростанням, наприклад 2 3 3 5 7 6 6 4 3 3. 

«Власна післягора» – це гора, що завершується спаданням, наприклад  3 5 7 7 6 6 4 3.

«Власна гора» – це гора, що починається зростанням, а завершується спаданням.

«Долина» – послідовність чисел, що йдуть поспіль і спочатку не зростають, а потім не спадають, наприклад 7 7 6 6 3 3 4 5 5 6 6.

«Власна переддолина» – це долина, що починається спаданням, наприклад 7 6 6 3 3 4 5 5 6 6. 

«Власна післядолина» – це долина, що завершується зростанням, наприклад 6 6 3 3 4 5 5 6.

\par \bigskip
\newpage

{\center  Варіанти \par \bigskip\bigskip}

%begin
1) \space 
З послідовності цілих вилучили всі числа, у сімковому запису яких нема непарних цифр. Знайти та вивести першу найдовшу підпослідовність залишку, що є парною власною пилкою та має довжину не менше двох.
%парною власною пилкою
%у сімковому запису яких нема непарних цифр

\par
%end
\vspace{0.9cm}
%begin
2) \space 
З послідовності цілих вилучили всі числа, у трійковому запису яких нема парних цифр. Знайти та вивести першу найдовшу підпослідовність залишку, що є власною серією та має довжину не менше двох.
%власною серією
%у трійковому запису яких нема парних цифр

\par
%end
\vspace{0.9cm}
%begin
3) \space 
З послідовності цілих вилучили всі числа, у сімковому запису яких нема цифри 2. Знайти та вивести першу найдовшу підпослідовність залишку, що є парною пилкою та має довжину не менше двох.
%парною пилкою
%у сімковому запису яких нема цифри 2

\par
%end
\vspace{0.9cm}
%begin
4) \space 
З послідовності цілих вилучили всі числа, у трійковому запису яких присутня цифра 2. Знайти та вивести першу найдовшу підпослідовність залишку, що є власною горою та має довжину не менше двох.
%власною горою
%у трійковому запису яких присутня цифра 2

\par
%end
\vspace{0.9cm}
%begin
5) \space 
З послідовності цілих вилучили всі числа, у сімковому запису яких відсутня цифра 5. Знайти та вивести першу найдовшу підпослідовність залишку, що є власною зростаючою серією та має довжину не менше двох.
%власною зростаючою серією
%у сімковому запису яких відсутня цифра 5

\par
%end
\vspace{0.9cm}
%begin
6) \space 
З послідовності цілих вилучили всі числа, які мають більше чотирьох різних натуральних дільників. Знайти та вивести першу найдовшу підпослідовність залишку, що є піком та має довжину не менше двох.
%піком
%які мають більше чотирьох різних натуральних дільників

\par
%end
\vspace{0.9cm}
%begin
7) \space 
З послідовності цілих вилучили всі числа, у сімковому запису яких присутня цифра 5. Знайти та вивести першу найдовшу підпослідовність залишку, що є долиною та має довжину не менше двох.
%долиною
%у сімковому запису яких присутня цифра 5

\par
%end
\vspace{0.9cm}
%begin
8) \space 
З послідовності цілих вилучили всі числа, які мають не більше п'яти різних натуральних дільників. Знайти та вивести першу найдовшу підпослідовність залишку, що є зростаючою серією та має довжину не менше двох.
%зростаючою серією
%які мають не більше п'яти різних натуральних дільників

\par
%end
\vspace{0.9cm}
%begin
9) \space 
З послідовності цілих вилучили всі числа, у сімковому запису яких нема непарних цифр. Знайти та вивести першу найдовшу підпослідовність залишку, що є власною пилкою та має довжину не менше двох.
%власною пилкою
%у сімковому запису яких нема непарних цифр

\par
%end
\vspace{0.9cm}
%begin
10) \space 
З послідовності цілих вилучили всі числа, у трійковому запису яких присутня цифра 2. Знайти та вивести першу найдовшу підпослідовність залишку, що є власною спадною серією та має довжину не менше двох.
%власною спадною серією
%у трійковому запису яких присутня цифра 2

\par
%end
\vspace{0.9cm}
%begin
11) \space 
З послідовності цілих вилучили всі числа, у сімковому запису яких відсутня цифра 5. Знайти та вивести першу найдовшу підпослідовність залишку, що є горою та має довжину не менше двох.
%горою
%у сімковому запису яких відсутня цифра 5

\par
%end
\vspace{0.9cm}
%begin
12) \space 
З послідовності цілих вилучили всі числа, які мають не більше п'яти різних натуральних дільників. Знайти та вивести першу найдовшу підпослідовність залишку, що є власною післягорою та має довжину не менше двох.
%власною післягорою
%які мають не більше п'яти різних натуральних дільників

\par
%end
\vspace{0.9cm}
%begin
13) \space 
З послідовності цілих вилучили всі числа, які мають більше чотирьох різних натуральних дільників. Знайти та вивести першу найдовшу підпослідовність залишку, що є власною передгорою та має довжину не менше двох.
%власною передгорою
%які мають більше чотирьох різних натуральних дільників

\par
%end
\vspace{0.9cm}
%begin
14) \space 
З послідовності цілих вилучили всі числа, у трійковому запису яких нема парних цифр. Знайти та вивести першу найдовшу підпослідовність залишку, що є ямою та має довжину не менше двох.
%ямою
%у трійковому запису яких нема парних цифр

\par
%end
\vspace{0.9cm}
%begin
15) \space 
З послідовності цілих вилучили всі числа, у сімковому запису яких присутня цифра 5. Знайти та вивести першу найдовшу підпослідовність залишку, що є непарною власною пилкою та має довжину не менше двох.
%непарною власною пилкою
%у сімковому запису яких присутня цифра 5

\par
%end
\vspace{0.9cm}
%begin
16) \space 
З послідовності цілих вилучили всі числа, у сімковому запису яких нема цифри 2. Знайти та вивести першу найдовшу підпослідовність залишку, що є власною переддолиною та має довжину не менше двох.
%власною переддолиною
%у сімковому запису яких нема цифри 2

\par
%end
\vspace{0.9cm}
%begin
17) \space 
З послідовності цілих вилучили всі числа, які мають не більше п'яти різних натуральних дільників. Знайти та вивести першу найдовшу підпослідовність залишку, що є власною післядолиною та має довжину не менше двох.
%власною післядолиною
%які мають не більше п'яти різних натуральних дільників

\par
%end
\vspace{0.9cm}
%begin
18) \space 
З послідовності цілих вилучили всі числа, які мають більше чотирьох різних натуральних дільників. Знайти та вивести першу найдовшу підпослідовність залишку, що є воронкою та має довжину не менше двох.
%воронкою
%які мають більше чотирьох різних натуральних дільників

\par
%end
\vspace{0.9cm}
%begin
19) \space 
З послідовності цілих вилучили всі числа, у трійковому запису яких нема парних цифр. Знайти та вивести першу найдовшу підпослідовність залишку, що є непарною пилкою та має довжину не менше двох.
%непарною пилкою
%у трійковому запису яких нема парних цифр

\par
%end
\vspace{0.9cm}
%begin
20) \space 
З послідовності цілих вилучили всі числа, у трійковому запису яких присутня цифра 2. Знайти та вивести першу найдовшу підпослідовність залишку, що є спадною серією та має довжину не менше двох.
%спадною серією
%у трійковому запису яких присутня цифра 2

\par
%end
\vspace{0.9cm}
%begin
21) \space 
З послідовності цілих вилучили всі числа, у сімковому запису яких присутня цифра 5. Знайти та вивести першу найдовшу підпослідовність залишку, що є плато та має довжину не менше двох.
%плато
%у сімковому запису яких присутня цифра 5

\par
%end
\vspace{0.9cm}
%begin
22) \space 
З послідовності цілих вилучили всі числа, у сімковому запису яких відсутня цифра 5. Знайти та вивести першу найдовшу підпослідовність залишку, що є серією та має довжину не менше двох.
%серією
%у сімковому запису яких відсутня цифра 5

\par
%end
\vspace{0.9cm}
%begin
23) \space 
З послідовності цілих вилучили всі числа, у сімковому запису яких нема цифри 2. Знайти та вивести першу найдовшу підпослідовність залишку, що є власною пилкою та має довжину не менше двох.
%власною пилкою
%у сімковому запису яких нема цифри 2

\par
%end
\vspace{0.9cm}
%begin
24) \space 
З послідовності цілих вилучили всі числа, у сімковому запису яких нема непарних цифр. Знайти та вивести першу найдовшу підпослідовність залишку, що є парною пилкою та має довжину не менше двох.
%парною пилкою
%у сімковому запису яких нема непарних цифр

\par
%end
\vspace{0.9cm}
%begin
25) \space 
З послідовності цілих вилучили всі числа, у трійковому запису яких присутня цифра 2. Знайти та вивести першу найдовшу підпослідовність залишку, що є воронкою та має довжину не менше двох.
%воронкою
%у трійковому запису яких присутня цифра 2

\par
%end
\vspace{0.9cm}
%begin
26) \space 
З послідовності цілих вилучили всі числа, у сімковому запису яких відсутня цифра 5. Знайти та вивести першу найдовшу підпослідовність залишку, що є власною горою та має довжину не менше двох.
%власною горою
%у сімковому запису яких відсутня цифра 5

\par
%end
\vspace{0.9cm}
%begin
27) \space 
З послідовності цілих вилучили всі числа, у сімковому запису яких нема цифри 2. Знайти та вивести першу найдовшу підпослідовність залишку, що є долиною та має довжину не менше двох.
%долиною
%у сімковому запису яких нема цифри 2

\par
%end
\vspace{0.9cm}
%begin
28) \space 
З послідовності цілих вилучили всі числа, у сімковому запису яких нема непарних цифр. Знайти та вивести першу найдовшу підпослідовність залишку, що є непарною власною пилкою та має довжину не менше двох.
%непарною власною пилкою
%у сімковому запису яких нема непарних цифр

\par
%end
\vspace{0.9cm}
%begin
29) \space 
З послідовності цілих вилучили всі числа, які мають більше чотирьох різних натуральних дільників. Знайти та вивести першу найдовшу підпослідовність залишку, що є власною спадною серією та має довжину не менше двох.
%власною спадною серією
%які мають більше чотирьох різних натуральних дільників

\par
%end
\vspace{0.9cm}
%begin
30) \space 
З послідовності цілих вилучили всі числа, у сімковому запису яких присутня цифра 5. Знайти та вивести першу найдовшу підпослідовність залишку, що є серією та має довжину не менше двох.
%серією
%у сімковому запису яких присутня цифра 5

\par
%end
\vspace{0.9cm}
%begin
31) \space 
З послідовності цілих вилучили всі числа, які мають не більше п'яти різних натуральних дільників. Знайти та вивести першу найдовшу підпослідовність залишку, що є власною післягорою та має довжину не менше двох.
%власною післягорою
%які мають не більше п'яти різних натуральних дільників

\par
%end
\vspace{0.9cm}
%begin
32) \space 
З послідовності цілих вилучили всі числа, у трійковому запису яких нема парних цифр. Знайти та вивести першу найдовшу підпослідовність залишку, що є піком та має довжину не менше двох.
%піком
%у трійковому запису яких нема парних цифр

\par
%end
\vspace{0.9cm}
%begin
33) \space 
З послідовності цілих вилучили всі числа, у сімковому запису яких нема непарних цифр. Знайти та вивести першу найдовшу підпослідовність залишку, що є власною зростаючою серією та має довжину не менше двох.
%власною зростаючою серією
%у сімковому запису яких нема непарних цифр

\par
%end
\vspace{0.9cm}
%begin
34) \space 
З послідовності цілих вилучили всі числа, які мають більше чотирьох різних натуральних дільників. Знайти та вивести першу найдовшу підпослідовність залишку, що є спадною серією та має довжину не менше двох.
%спадною серією
%які мають більше чотирьох різних натуральних дільників

\par
%end
\vspace{0.9cm}
%begin
35) \space 
З послідовності цілих вилучили всі числа, у сімковому запису яких нема цифри 2. Знайти та вивести першу найдовшу підпослідовність залишку, що є ямою та має довжину не менше двох.
%ямою
%у сімковому запису яких нема цифри 2

\par
%end
\vspace{0.9cm}
%begin
36) \space 
З послідовності цілих вилучили всі числа, у трійковому запису яких присутня цифра 2. Знайти та вивести першу найдовшу підпослідовність залишку, що є власною післядолиною та має довжину не менше двох.
%власною післядолиною
%у трійковому запису яких присутня цифра 2

\par
%end
\vspace{0.9cm}
%begin
37) \space 
З послідовності цілих вилучили всі числа, у сімковому запису яких відсутня цифра 5. Знайти та вивести першу найдовшу підпослідовність залишку, що є парною власною пилкою та має довжину не менше двох.
%парною власною пилкою
%у сімковому запису яких відсутня цифра 5

\par
%end
\vspace{0.9cm}
%begin
38) \space 
З послідовності цілих вилучили всі числа, у трійковому запису яких нема парних цифр. Знайти та вивести першу найдовшу підпослідовність залишку, що є зростаючою серією та має довжину не менше двох.
%зростаючою серією
%у трійковому запису яких нема парних цифр

\par
%end
\vspace{0.9cm}
%begin
39) \space 
З послідовності цілих вилучили всі числа, у сімковому запису яких присутня цифра 5. Знайти та вивести першу найдовшу підпослідовність залишку, що є плато та має довжину не менше двох.
%плато
%у сімковому запису яких присутня цифра 5

\par
%end
\vspace{0.9cm}
%begin
40) \space 
З послідовності цілих вилучили всі числа, які мають не більше п'яти різних натуральних дільників. Знайти та вивести першу найдовшу підпослідовність залишку, що є власною передгорою та має довжину не менше двох.
%власною передгорою
%які мають не більше п'яти різних натуральних дільників

\par
%end
\vspace{0.9cm}
%begin
41) \space 
З послідовності цілих вилучили всі числа, у сімковому запису яких нема цифри 2. Знайти та вивести першу найдовшу підпослідовність залишку, що є горою та має довжину не менше двох.
%горою
%у сімковому запису яких нема цифри 2

\par
%end
\vspace{0.9cm}
%begin
42) \space 
З послідовності цілих вилучили всі числа, у трійковому запису яких нема парних цифр. Знайти та вивести першу найдовшу підпослідовність залишку, що є непарною пилкою та має довжину не менше двох.
%непарною пилкою
%у трійковому запису яких нема парних цифр

\par
%end
\vspace{0.9cm}
%begin
43) \space 
З послідовності цілих вилучили всі числа, у сімковому запису яких присутня цифра 5. Знайти та вивести першу найдовшу підпослідовність залишку, що є власною переддолиною та має довжину не менше двох.
%власною переддолиною
%у сімковому запису яких присутня цифра 5

\par
%end
\vspace{0.9cm}
%begin
44) \space 
З послідовності цілих вилучили всі числа, які мають більше чотирьох різних натуральних дільників. Знайти та вивести першу найдовшу підпослідовність залишку, що є власною серією та має довжину не менше двох.
%власною серією
%які мають більше чотирьох різних натуральних дільників

\par
%end
\vspace{0.9cm}
%begin
45) \space 
З послідовності цілих вилучили всі числа, які мають не більше п'яти різних натуральних дільників. Знайти та вивести першу найдовшу підпослідовність залишку, що є спадною серією та має довжину не менше двох.
%спадною серією
%які мають не більше п'яти різних натуральних дільників

\par
%end
\vspace{0.9cm}
%begin
46) \space 
З послідовності цілих вилучили всі числа, у сімковому запису яких нема непарних цифр. Знайти та вивести першу найдовшу підпослідовність залишку, що є власною післягорою та має довжину не менше двох.
%власною післягорою
%у сімковому запису яких нема непарних цифр

\par
%end
\vspace{0.9cm}
%begin
47) \space 
З послідовності цілих вилучили всі числа, у сімковому запису яких відсутня цифра 5. Знайти та вивести першу найдовшу підпослідовність залишку, що є долиною та має довжину не менше двох.
%долиною
%у сімковому запису яких відсутня цифра 5

\par
%end
\vspace{0.9cm}
%begin
48) \space 
З послідовності цілих вилучили всі числа, у трійковому запису яких присутня цифра 2. Знайти та вивести першу найдовшу підпослідовність залишку, що є горою та має довжину не менше двох.
%горою
%у трійковому запису яких присутня цифра 2

\par
%end
\vspace{0.9cm}
%begin
49) \space 
З послідовності цілих вилучили всі числа, у сімковому запису яких нема непарних цифр. Знайти та вивести першу найдовшу підпослідовність залишку, що є власною серією та має довжину не менше двох.
%власною серією
%у сімковому запису яких нема непарних цифр

\par
%end
\vspace{0.9cm}
%begin
50) \space 
З послідовності цілих вилучили всі числа, у сімковому запису яких відсутня цифра 5. Знайти та вивести першу найдовшу підпослідовність залишку, що є парною власною пилкою та має довжину не менше двох.
%парною власною пилкою
%у сімковому запису яких відсутня цифра 5

\par
%end
\vspace{0.9cm}
%begin
51) \space 
З послідовності цілих вилучили всі числа, які мають не більше п'яти різних натуральних дільників. Знайти та вивести першу найдовшу підпослідовність залишку, що є парною пилкою та має довжину не менше двох.
%парною пилкою
%які мають не більше п'яти різних натуральних дільників

\par
%end
\vspace{0.9cm}
%begin
52) \space 
З послідовності цілих вилучили всі числа, у трійковому запису яких нема парних цифр. Знайти та вивести першу найдовшу підпослідовність залишку, що є власною зростаючою серією та має довжину не менше двох.
%власною зростаючою серією
%у трійковому запису яких нема парних цифр

\par
%end
\vspace{0.9cm}
%begin
53) \space 
З послідовності цілих вилучили всі числа, у сімковому запису яких присутня цифра 5. Знайти та вивести першу найдовшу підпослідовність залишку, що є власною спадною серією та має довжину не менше двох.
%власною спадною серією
%у сімковому запису яких присутня цифра 5

\par
%end
\vspace{0.9cm}
%begin
54) \space 
З послідовності цілих вилучили всі числа, у трійковому запису яких присутня цифра 2. Знайти та вивести першу найдовшу підпослідовність залишку, що є плато та має довжину не менше двох.
%плато
%у трійковому запису яких присутня цифра 2

\par
%end
\vspace{0.9cm}
%begin
55) \space 
З послідовності цілих вилучили всі числа, у сімковому запису яких нема цифри 2. Знайти та вивести першу найдовшу підпослідовність залишку, що є власною передгорою та має довжину не менше двох.
%власною передгорою
%у сімковому запису яких нема цифри 2

\par
%end
\vspace{0.9cm}
%begin
56) \space 
З послідовності цілих вилучили всі числа, які мають більше чотирьох різних натуральних дільників. Знайти та вивести першу найдовшу підпослідовність залишку, що є власною післядолиною та має довжину не менше двох.
%власною післядолиною
%які мають більше чотирьох різних натуральних дільників

\par
%end
\vspace{0.9cm}
%begin
57) \space 
З послідовності цілих вилучили всі числа, у трійковому запису яких присутня цифра 2. Знайти та вивести першу найдовшу підпослідовність залишку, що є ямою та має довжину не менше двох.
%ямою
%у трійковому запису яких присутня цифра 2

\par
%end
\vspace{0.9cm}
%begin
58) \space 
З послідовності цілих вилучили всі числа, у сімковому запису яких присутня цифра 5. Знайти та вивести першу найдовшу підпослідовність залишку, що є воронкою та має довжину не менше двох.
%воронкою
%у сімковому запису яких присутня цифра 5

\par
%end
\vspace{0.9cm}
%begin
59) \space 
З послідовності цілих вилучили всі числа, у трійковому запису яких нема парних цифр. Знайти та вивести першу найдовшу підпослідовність залишку, що є непарною пилкою та має довжину не менше двох.
%непарною пилкою
%у трійковому запису яких нема парних цифр

\par
%end
\vspace{0.9cm}
%begin
60) \space 
З послідовності цілих вилучили всі числа, у сімковому запису яких відсутня цифра 5. Знайти та вивести першу найдовшу підпослідовність залишку, що є непарною власною пилкою та має довжину не менше двох.
%непарною власною пилкою
%у сімковому запису яких відсутня цифра 5

\par
%end
\vspace{0.9cm}
%begin
61) \space 
З послідовності цілих вилучили всі числа, у сімковому запису яких нема цифри 2. Знайти та вивести першу найдовшу підпослідовність залишку, що є серією та має довжину не менше двох.
%серією
%у сімковому запису яких нема цифри 2

\par
%end
\vspace{0.9cm}
%begin
62) \space 
З послідовності цілих вилучили всі числа, у сімковому запису яких нема непарних цифр. Знайти та вивести першу найдовшу підпослідовність залишку, що є зростаючою серією та має довжину не менше двох.
%зростаючою серією
%у сімковому запису яких нема непарних цифр

\par
%end
\vspace{0.9cm}
%begin
63) \space 
З послідовності цілих вилучили всі числа, які мають більше чотирьох різних натуральних дільників. Знайти та вивести першу найдовшу підпослідовність залишку, що є власною пилкою та має довжину не менше двох.
%власною пилкою
%які мають більше чотирьох різних натуральних дільників

\par
%end
\vspace{0.9cm}
%begin
64) \space 
З послідовності цілих вилучили всі числа, які мають не більше п'яти різних натуральних дільників. Знайти та вивести першу найдовшу підпослідовність залишку, що є власною горою та має довжину не менше двох.
%власною горою
%які мають не більше п'яти різних натуральних дільників

\par
%end
\vspace{0.9cm}
%begin
65) \space 
З послідовності цілих вилучили всі числа, у сімковому запису яких нема цифри 2. Знайти та вивести першу найдовшу підпослідовність залишку, що є власною переддолиною та має довжину не менше двох.
%власною переддолиною
%у сімковому запису яких нема цифри 2

\par
%end
\vspace{0.9cm}
%begin
66) \space 
З послідовності цілих вилучили всі числа, у сімковому запису яких присутня цифра 5. Знайти та вивести першу найдовшу підпослідовність залишку, що є піком та має довжину не менше двох.
%піком
%у сімковому запису яких присутня цифра 5

\par
%end
\vspace{0.9cm}
%begin
67) \space 
З послідовності цілих вилучили всі числа, які мають не більше п'яти різних натуральних дільників. Знайти та вивести першу найдовшу підпослідовність залишку, що є власною серією та має довжину не менше двох.
%власною серією
%які мають не більше п'яти різних натуральних дільників

\par
%end
\vspace{0.9cm}
%begin
68) \space 
З послідовності цілих вилучили всі числа, які мають більше чотирьох різних натуральних дільників. Знайти та вивести першу найдовшу підпослідовність залишку, що є долиною та має довжину не менше двох.
%долиною
%які мають більше чотирьох різних натуральних дільників

\par
%end
\vspace{0.9cm}
%begin
69) \space 
З послідовності цілих вилучили всі числа, у трійковому запису яких нема парних цифр. Знайти та вивести першу найдовшу підпослідовність залишку, що є воронкою та має довжину не менше двох.
%воронкою
%у трійковому запису яких нема парних цифр

\par
%end
\vspace{0.9cm}
%begin
70) \space 
З послідовності цілих вилучили всі числа, у трійковому запису яких присутня цифра 2. Знайти та вивести першу найдовшу підпослідовність залишку, що є спадною серією та має довжину не менше двох.
%спадною серією
%у трійковому запису яких присутня цифра 2

\par
%end
\vspace{0.9cm}
%begin
71) \space 
З послідовності цілих вилучили всі числа, у сімковому запису яких нема непарних цифр. Знайти та вивести першу найдовшу підпослідовність залишку, що є парною власною пилкою та має довжину не менше двох.
%парною власною пилкою
%у сімковому запису яких нема непарних цифр

\par
%end
\vspace{0.9cm}
%begin
72) \space 
З послідовності цілих вилучили всі числа, у сімковому запису яких відсутня цифра 5. Знайти та вивести першу найдовшу підпослідовність залишку, що є власною післягорою та має довжину не менше двох.
%власною післягорою
%у сімковому запису яких відсутня цифра 5

\par
%end
\vspace{0.9cm}
%begin
73) \space 
З послідовності цілих вилучили всі числа, у сімковому запису яких нема цифри 2. Знайти та вивести першу найдовшу підпослідовність залишку, що є власною передгорою та має довжину не менше двох.
%власною передгорою
%у сімковому запису яких нема цифри 2

\par
%end
\vspace{0.9cm}
%begin
74) \space 
З послідовності цілих вилучили всі числа, які мають не більше п'яти різних натуральних дільників. Знайти та вивести першу найдовшу підпослідовність залишку, що є власною спадною серією та має довжину не менше двох.
%власною спадною серією
%які мають не більше п'яти різних натуральних дільників

\par
%end
\vspace{0.9cm}
%begin
75) \space 
З послідовності цілих вилучили всі числа, у трійковому запису яких присутня цифра 2. Знайти та вивести першу найдовшу підпослідовність залишку, що є власною горою та має довжину не менше двох.
%власною горою
%у трійковому запису яких присутня цифра 2

\par
%end
\vspace{0.9cm}
%begin
76) \space 
З послідовності цілих вилучили всі числа, у сімковому запису яких нема непарних цифр. Знайти та вивести першу найдовшу підпослідовність залишку, що є плато та має довжину не менше двох.
%плато
%у сімковому запису яких нема непарних цифр

\par
%end
\vspace{0.9cm}
%begin
77) \space 
З послідовності цілих вилучили всі числа, які мають більше чотирьох різних натуральних дільників. Знайти та вивести першу найдовшу підпослідовність залишку, що є горою та має довжину не менше двох.
%горою
%які мають більше чотирьох різних натуральних дільників

\par
%end
\vspace{0.9cm}
%begin
78) \space 
З послідовності цілих вилучили всі числа, у сімковому запису яких відсутня цифра 5. Знайти та вивести першу найдовшу підпослідовність залишку, що є непарною власною пилкою та має довжину не менше двох.
%непарною власною пилкою
%у сімковому запису яких відсутня цифра 5

\par
%end
\vspace{0.9cm}
%begin
79) \space 
З послідовності цілих вилучили всі числа, у сімковому запису яких присутня цифра 5. Знайти та вивести першу найдовшу підпослідовність залишку, що є власною післядолиною та має довжину не менше двох.
%власною післядолиною
%у сімковому запису яких присутня цифра 5

\par
%end
\vspace{0.9cm}
%begin
80) \space 
З послідовності цілих вилучили всі числа, у трійковому запису яких нема парних цифр. Знайти та вивести першу найдовшу підпослідовність залишку, що є непарною пилкою та має довжину не менше двох.
%непарною пилкою
%у трійковому запису яких нема парних цифр

\par
%end
\vspace{0.9cm}
%begin
81) \space 
З послідовності цілих вилучили всі числа, які мають більше чотирьох різних натуральних дільників. Знайти та вивести першу найдовшу підпослідовність залишку, що є серією та має довжину не менше двох.
%серією
%які мають більше чотирьох різних натуральних дільників

\par
%end
\vspace{0.9cm}
%begin
82) \space 
З послідовності цілих вилучили всі числа, які мають не більше п'яти різних натуральних дільників. Знайти та вивести першу найдовшу підпослідовність залишку, що є піком та має довжину не менше двох.
%піком
%які мають не більше п'яти різних натуральних дільників

\par
%end
\vspace{0.9cm}
%begin
83) \space 
З послідовності цілих вилучили всі числа, у сімковому запису яких відсутня цифра 5. Знайти та вивести першу найдовшу підпослідовність залишку, що є власною переддолиною та має довжину не менше двох.
%власною переддолиною
%у сімковому запису яких відсутня цифра 5

\par
%end
\vspace{0.9cm}
%begin
84) \space 
З послідовності цілих вилучили всі числа, у сімковому запису яких присутня цифра 5. Знайти та вивести першу найдовшу підпослідовність залишку, що є зростаючою серією та має довжину не менше двох.
%зростаючою серією
%у сімковому запису яких присутня цифра 5

\par
%end
\vspace{0.9cm}
%begin
85) \space 
З послідовності цілих вилучили всі числа, у сімковому запису яких нема непарних цифр. Знайти та вивести першу найдовшу підпослідовність залишку, що є парною пилкою та має довжину не менше двох.
%парною пилкою
%у сімковому запису яких нема непарних цифр

\par
%end
\vspace{0.9cm}
%begin
86) \space 
З послідовності цілих вилучили всі числа, у трійковому запису яких присутня цифра 2. Знайти та вивести першу найдовшу підпослідовність залишку, що є ямою та має довжину не менше двох.
%ямою
%у трійковому запису яких присутня цифра 2

\par
%end
\vspace{0.9cm}
%begin
87) \space 
З послідовності цілих вилучили всі числа, у трійковому запису яких нема парних цифр. Знайти та вивести першу найдовшу підпослідовність залишку, що є власною зростаючою серією та має довжину не менше двох.
%власною зростаючою серією
%у трійковому запису яких нема парних цифр

\par
%end
\vspace{0.9cm}
%begin
88) \space 
З послідовності цілих вилучили всі числа, у сімковому запису яких нема цифри 2. Знайти та вивести першу найдовшу підпослідовність залишку, що є власною пилкою та має довжину не менше двох.
%власною пилкою
%у сімковому запису яких нема цифри 2

\par
%end
\vspace{0.9cm}
%begin
89) \space 
З послідовності цілих вилучили всі числа, які мають не більше п'яти різних натуральних дільників. Знайти та вивести першу найдовшу підпослідовність залишку, що є плато та має довжину не менше двох.
%плато
%які мають не більше п'яти різних натуральних дільників

\par
%end
\vspace{0.9cm}
%begin
90) \space 
З послідовності цілих вилучили всі числа, у сімковому запису яких присутня цифра 5. Знайти та вивести першу найдовшу підпослідовність залишку, що є непарною пилкою та має довжину не менше двох.
%непарною пилкою
%у сімковому запису яких присутня цифра 5

\par
%end
\vspace{0.9cm}
%begin
91) \space 
З послідовності цілих вилучили всі числа, у трійковому запису яких присутня цифра 2. Знайти та вивести першу найдовшу підпослідовність залишку, що є долиною та має довжину не менше двох.
%долиною
%у трійковому запису яких присутня цифра 2

\par
%end
\vspace{0.9cm}
%begin
92) \space 
З послідовності цілих вилучили всі числа, у трійковому запису яких нема парних цифр. Знайти та вивести першу найдовшу підпослідовність залишку, що є парною пилкою та має довжину не менше двох.
%парною пилкою
%у трійковому запису яких нема парних цифр

\par
%end
\vspace{0.9cm}
%begin
93) \space 
З послідовності цілих вилучили всі числа, у сімковому запису яких нема цифри 2. Знайти та вивести першу найдовшу підпослідовність залишку, що є власною серією та має довжину не менше двох.
%власною серією
%у сімковому запису яких нема цифри 2

\par
%end
\vspace{0.9cm}
%begin
94) \space 
З послідовності цілих вилучили всі числа, у сімковому запису яких відсутня цифра 5. Знайти та вивести першу найдовшу підпослідовність залишку, що є спадною серією та має довжину не менше двох.
%спадною серією
%у сімковому запису яких відсутня цифра 5

\par
%end
\vspace{0.9cm}
%begin
95) \space 
З послідовності цілих вилучили всі числа, у сімковому запису яких нема непарних цифр. Знайти та вивести першу найдовшу підпослідовність залишку, що є власною пилкою та має довжину не менше двох.
%власною пилкою
%у сімковому запису яких нема непарних цифр

\par
%end
\vspace{0.9cm}
%begin
96) \space 
З послідовності цілих вилучили всі числа, які мають більше чотирьох різних натуральних дільників. Знайти та вивести першу найдовшу підпослідовність залишку, що є власною післядолиною та має довжину не менше двох.
%власною післядолиною
%які мають більше чотирьох різних натуральних дільників

\par
%end
\vspace{0.9cm}
%begin
97) \space 
З послідовності цілих вилучили всі числа, у сімковому запису яких відсутня цифра 5. Знайти та вивести першу найдовшу підпослідовність залишку, що є серією та має довжину не менше двох.
%серією
%у сімковому запису яких відсутня цифра 5

\par
%end
\vspace{0.9cm}
%begin
98) \space 
З послідовності цілих вилучили всі числа, у сімковому запису яких нема непарних цифр. Знайти та вивести першу найдовшу підпослідовність залишку, що є воронкою та має довжину не менше двох.
%воронкою
%у сімковому запису яких нема непарних цифр

\par
%end
\vspace{0.9cm}
%begin
99) \space 
З послідовності цілих вилучили всі числа, у трійковому запису яких присутня цифра 2. Знайти та вивести першу найдовшу підпослідовність залишку, що є непарною власною пилкою та має довжину не менше двох.
%непарною власною пилкою
%у трійковому запису яких присутня цифра 2

\par
%end
\vspace{0.9cm}
%begin
100) \space 
З послідовності цілих вилучили всі числа, у трійковому запису яких нема парних цифр. Знайти та вивести першу найдовшу підпослідовність залишку, що є парною власною пилкою та має довжину не менше двох.
%парною власною пилкою
%у трійковому запису яких нема парних цифр

\par
%end
\vspace{0.9cm}
%begin
101) \space 
З послідовності цілих вилучили всі числа, у сімковому запису яких присутня цифра 5. Знайти та вивести першу найдовшу підпослідовність залишку, що є зростаючою серією та має довжину не менше двох.
%зростаючою серією
%у сімковому запису яких присутня цифра 5

\par
%end
\vspace{0.9cm}
%begin
102) \space 
З послідовності цілих вилучили всі числа, які мають не більше п'яти різних натуральних дільників. Знайти та вивести першу найдовшу підпослідовність залишку, що є власною спадною серією та має довжину не менше двох.
%власною спадною серією
%які мають не більше п'яти різних натуральних дільників

\par
%end
\vspace{0.9cm}
%begin
103) \space 
З послідовності цілих вилучили всі числа, у сімковому запису яких нема цифри 2. Знайти та вивести першу найдовшу підпослідовність залишку, що є власною горою та має довжину не менше двох.
%власною горою
%у сімковому запису яких нема цифри 2

\par
%end
\vspace{0.9cm}
%begin
104) \space 
З послідовності цілих вилучили всі числа, які мають більше чотирьох різних натуральних дільників. Знайти та вивести першу найдовшу підпослідовність залишку, що є власною післягорою та має довжину не менше двох.
%власною післягорою
%які мають більше чотирьох різних натуральних дільників

\par
%end
\vspace{0.9cm}
%begin
105) \space 
З послідовності цілих вилучили всі числа, у сімковому запису яких присутня цифра 5. Знайти та вивести першу найдовшу підпослідовність залишку, що є власною передгорою та має довжину не менше двох.
%власною передгорою
%у сімковому запису яких присутня цифра 5

\par
%end
\vspace{0.9cm}
%begin
106) \space 
З послідовності цілих вилучили всі числа, які мають не більше п'яти різних натуральних дільників. Знайти та вивести першу найдовшу підпослідовність залишку, що є піком та має довжину не менше двох.
%піком
%які мають не більше п'яти різних натуральних дільників

\par
%end
\vspace{0.9cm}
%begin
107) \space 
З послідовності цілих вилучили всі числа, які мають більше чотирьох різних натуральних дільників. Знайти та вивести першу найдовшу підпослідовність залишку, що є власною зростаючою серією та має довжину не менше двох.
%власною зростаючою серією
%які мають більше чотирьох різних натуральних дільників

\par
%end
\vspace{0.9cm}
%begin
108) \space 
З послідовності цілих вилучили всі числа, у сімковому запису яких нема цифри 2. Знайти та вивести першу найдовшу підпослідовність залишку, що є горою та має довжину не менше двох.
%горою
%у сімковому запису яких нема цифри 2

\par
%end
\vspace{0.9cm}
%begin
109) \space 
З послідовності цілих вилучили всі числа, у сімковому запису яких відсутня цифра 5. Знайти та вивести першу найдовшу підпослідовність залишку, що є ямою та має довжину не менше двох.
%ямою
%у сімковому запису яких відсутня цифра 5

\par
%end
\vspace{0.9cm}
%begin
110) \space 
З послідовності цілих вилучили всі числа, у трійковому запису яких нема парних цифр. Знайти та вивести першу найдовшу підпослідовність залишку, що є власною переддолиною та має довжину не менше двох.
%власною переддолиною
%у трійковому запису яких нема парних цифр

\par
%end
\vspace{0.9cm}
%begin
111) \space 
З послідовності цілих вилучили всі числа, у сімковому запису яких нема непарних цифр. Знайти та вивести першу найдовшу підпослідовність залишку, що є власною переддолиною та має довжину не менше двох.
%власною переддолиною
%у сімковому запису яких нема непарних цифр

\par
%end
\vspace{0.9cm}
%begin
112) \space 
З послідовності цілих вилучили всі числа, у трійковому запису яких присутня цифра 2. Знайти та вивести першу найдовшу підпослідовність залишку, що є серією та має довжину не менше двох.
%серією
%у трійковому запису яких присутня цифра 2

\par
%end
\vspace{0.9cm}
%begin
113) \space 
З послідовності цілих вилучили всі числа, які мають не більше п'яти різних натуральних дільників. Знайти та вивести першу найдовшу підпослідовність залишку, що є власною спадною серією та має довжину не менше двох.
%власною спадною серією
%які мають не більше п'яти різних натуральних дільників

\par
%end
\vspace{0.9cm}
%begin
114) \space 
З послідовності цілих вилучили всі числа, у трійковому запису яких присутня цифра 2. Знайти та вивести першу найдовшу підпослідовність залишку, що є непарною власною пилкою та має довжину не менше двох.
%непарною власною пилкою
%у трійковому запису яких присутня цифра 2

\par
%end
\vspace{0.9cm}
%begin
115) \space 
З послідовності цілих вилучили всі числа, у сімковому запису яких нема непарних цифр. Знайти та вивести першу найдовшу підпослідовність залишку, що є власною пилкою та має довжину не менше двох.
%власною пилкою
%у сімковому запису яких нема непарних цифр

\par
%end
\vspace{0.9cm}
%begin
116) \space 
З послідовності цілих вилучили всі числа, які мають більше чотирьох різних натуральних дільників. Знайти та вивести першу найдовшу підпослідовність залишку, що є воронкою та має довжину не менше двох.
%воронкою
%які мають більше чотирьох різних натуральних дільників

\par
%end
\vspace{0.9cm}
%begin
117) \space 
З послідовності цілих вилучили всі числа, у сімковому запису яких нема цифри 2. Знайти та вивести першу найдовшу підпослідовність залишку, що є спадною серією та має довжину не менше двох.
%спадною серією
%у сімковому запису яких нема цифри 2

\par
%end
\vspace{0.9cm}
%begin
118) \space 
З послідовності цілих вилучили всі числа, у сімковому запису яких присутня цифра 5. Знайти та вивести першу найдовшу підпослідовність залишку, що є власною післядолиною та має довжину не менше двох.
%власною післядолиною
%у сімковому запису яких присутня цифра 5

\par
%end
\vspace{0.9cm}
%begin
119) \space 
З послідовності цілих вилучили всі числа, у трійковому запису яких нема парних цифр. Знайти та вивести першу найдовшу підпослідовність залишку, що є піком та має довжину не менше двох.
%піком
%у трійковому запису яких нема парних цифр

\par
%end
\vspace{0.9cm}
%begin
120) \space 
З послідовності цілих вилучили всі числа, у сімковому запису яких відсутня цифра 5. Знайти та вивести першу найдовшу підпослідовність залишку, що є власною серією та має довжину не менше двох.
%власною серією
%у сімковому запису яких відсутня цифра 5

\par
%end
\vspace{0.9cm}
%begin
121) \space 
З послідовності цілих вилучили всі числа, у сімковому запису яких нема непарних цифр. Знайти та вивести першу найдовшу підпослідовність залишку, що є зростаючою серією та має довжину не менше двох.
%зростаючою серією
%у сімковому запису яких нема непарних цифр

\par
%end
\vspace{0.9cm}
%begin
122) \space 
З послідовності цілих вилучили всі числа, у сімковому запису яких присутня цифра 5. Знайти та вивести першу найдовшу підпослідовність залишку, що є долиною та має довжину не менше двох.
%долиною
%у сімковому запису яких присутня цифра 5

\par
%end
\vspace{0.9cm}
%begin
123) \space 
З послідовності цілих вилучили всі числа, у трійковому запису яких присутня цифра 2. Знайти та вивести першу найдовшу підпослідовність залишку, що є власною передгорою та має довжину не менше двох.
%власною передгорою
%у трійковому запису яких присутня цифра 2

\par
%end
\vspace{0.9cm}
%begin
124) \space 
З послідовності цілих вилучили всі числа, у сімковому запису яких нема цифри 2. Знайти та вивести першу найдовшу підпослідовність залишку, що є плато та має довжину не менше двох.
%плато
%у сімковому запису яких нема цифри 2

\par
%end
\vspace{0.9cm}
%begin
125) \space 
З послідовності цілих вилучили всі числа, у трійковому запису яких нема парних цифр. Знайти та вивести першу найдовшу підпослідовність залишку, що є власною горою та має довжину не менше двох.
%власною горою
%у трійковому запису яких нема парних цифр

\par
%end
\vspace{0.9cm}
%begin
126) \space 
З послідовності цілих вилучили всі числа, які мають більше чотирьох різних натуральних дільників. Знайти та вивести першу найдовшу підпослідовність залишку, що є непарною пилкою та має довжину не менше двох.
%непарною пилкою
%які мають більше чотирьох різних натуральних дільників

\par
%end
\vspace{0.9cm}
%begin
127) \space 
З послідовності цілих вилучили всі числа, які мають не більше п'яти різних натуральних дільників. Знайти та вивести першу найдовшу підпослідовність залишку, що є горою та має довжину не менше двох.
%горою
%які мають не більше п'яти різних натуральних дільників

\par
%end
\vspace{0.9cm}
%begin
128) \space 
З послідовності цілих вилучили всі числа, у сімковому запису яких відсутня цифра 5. Знайти та вивести першу найдовшу підпослідовність залишку, що є ямою та має довжину не менше двох.
%ямою
%у сімковому запису яких відсутня цифра 5

\par
%end
\vspace{0.9cm}
%begin
129) \space 
З послідовності цілих вилучили всі числа, які мають більше чотирьох різних натуральних дільників. Знайти та вивести першу найдовшу підпослідовність залишку, що є парною пилкою та має довжину не менше двох.
%парною пилкою
%які мають більше чотирьох різних натуральних дільників

\par
%end
\vspace{0.9cm}
%begin
130) \space 
З послідовності цілих вилучили всі числа, у трійковому запису яких присутня цифра 2. Знайти та вивести першу найдовшу підпослідовність залишку, що є власною зростаючою серією та має довжину не менше двох.
%власною зростаючою серією
%у трійковому запису яких присутня цифра 2

\par
%end
\vspace{0.9cm}
%begin
131) \space 
З послідовності цілих вилучили всі числа, у сімковому запису яких нема цифри 2. Знайти та вивести першу найдовшу підпослідовність залишку, що є парною власною пилкою та має довжину не менше двох.
%парною власною пилкою
%у сімковому запису яких нема цифри 2

\par
%end
\vspace{0.9cm}
%begin
132) \space 
З послідовності цілих вилучили всі числа, у сімковому запису яких нема непарних цифр. Знайти та вивести першу найдовшу підпослідовність залишку, що є власною післягорою та має довжину не менше двох.
%власною післягорою
%у сімковому запису яких нема непарних цифр

\par
%end
\vspace{0.9cm}
%begin
133) \space 
З послідовності цілих вилучили всі числа, у сімковому запису яких відсутня цифра 5. Знайти та вивести першу найдовшу підпослідовність залишку, що є непарною пилкою та має довжину не менше двох.
%непарною пилкою
%у сімковому запису яких відсутня цифра 5

\par
%end
\vspace{0.9cm}
%begin
134) \space 
З послідовності цілих вилучили всі числа, які мають не більше п'яти різних натуральних дільників. Знайти та вивести першу найдовшу підпослідовність залишку, що є плато та має довжину не менше двох.
%плато
%які мають не більше п'яти різних натуральних дільників

\par
%end
\vspace{0.9cm}
%begin
135) \space 
З послідовності цілих вилучили всі числа, у трійковому запису яких нема парних цифр. Знайти та вивести першу найдовшу підпослідовність залишку, що є власною пилкою та має довжину не менше двох.
%власною пилкою
%у трійковому запису яких нема парних цифр

\par
%end
\vspace{0.9cm}
%begin
136) \space 
З послідовності цілих вилучили всі числа, у сімковому запису яких присутня цифра 5. Знайти та вивести першу найдовшу підпослідовність залишку, що є власною спадною серією та має довжину не менше двох.
%власною спадною серією
%у сімковому запису яких присутня цифра 5

\par
%end
\vspace{0.9cm}
%begin
137) \space 
З послідовності цілих вилучили всі числа, які мають не більше п'яти різних натуральних дільників. Знайти та вивести першу найдовшу підпослідовність залишку, що є власною горою та має довжину не менше двох.
%власною горою
%які мають не більше п'яти різних натуральних дільників

\par
%end
\vspace{0.9cm}
%begin
138) \space 
З послідовності цілих вилучили всі числа, у сімковому запису яких присутня цифра 5. Знайти та вивести першу найдовшу підпослідовність залишку, що є парною пилкою та має довжину не менше двох.
%парною пилкою
%у сімковому запису яких присутня цифра 5

\par
%end
\vspace{0.9cm}
%begin
139) \space 
З послідовності цілих вилучили всі числа, у сімковому запису яких нема непарних цифр. Знайти та вивести першу найдовшу підпослідовність залишку, що є власною післядолиною та має довжину не менше двох.
%власною післядолиною
%у сімковому запису яких нема непарних цифр

\par
%end
\vspace{0.9cm}
%begin
140) \space 
З послідовності цілих вилучили всі числа, які мають більше чотирьох різних натуральних дільників. Знайти та вивести першу найдовшу підпослідовність залишку, що є власною серією та має довжину не менше двох.
%власною серією
%які мають більше чотирьох різних натуральних дільників

\par
%end
\vspace{0.9cm}
%begin
141) \space 
З послідовності цілих вилучили всі числа, у трійковому запису яких присутня цифра 2. Знайти та вивести першу найдовшу підпослідовність залишку, що є піком та має довжину не менше двох.
%піком
%у трійковому запису яких присутня цифра 2

\par
%end
\vspace{0.9cm}
%begin
142) \space 
З послідовності цілих вилучили всі числа, у сімковому запису яких нема цифри 2. Знайти та вивести першу найдовшу підпослідовність залишку, що є непарною власною пилкою та має довжину не менше двох.
%непарною власною пилкою
%у сімковому запису яких нема цифри 2

\par
%end
\vspace{0.9cm}
%begin
143) \space 
З послідовності цілих вилучили всі числа, у сімковому запису яких відсутня цифра 5. Знайти та вивести першу найдовшу підпослідовність залишку, що є ямою та має довжину не менше двох.
%ямою
%у сімковому запису яких відсутня цифра 5

\par
%end
\vspace{0.9cm}
%begin
144) \space 
З послідовності цілих вилучили всі числа, у трійковому запису яких нема парних цифр. Знайти та вивести першу найдовшу підпослідовність залишку, що є серією та має довжину не менше двох.
%серією
%у трійковому запису яких нема парних цифр

\par
%end
\vspace{0.9cm}
%begin
145) \space 
З послідовності цілих вилучили всі числа, у сімковому запису яких відсутня цифра 5. Знайти та вивести першу найдовшу підпослідовність залишку, що є власною переддолиною та має довжину не менше двох.
%власною переддолиною
%у сімковому запису яких відсутня цифра 5

\par
%end
\vspace{0.9cm}
%begin
146) \space 
З послідовності цілих вилучили всі числа, у трійковому запису яких присутня цифра 2. Знайти та вивести першу найдовшу підпослідовність залишку, що є горою та має довжину не менше двох.
%горою
%у трійковому запису яких присутня цифра 2

\par
%end
\vspace{0.9cm}
%begin
147) \space 
З послідовності цілих вилучили всі числа, які мають більше чотирьох різних натуральних дільників. Знайти та вивести першу найдовшу підпослідовність залишку, що є власною зростаючою серією та має довжину не менше двох.
%власною зростаючою серією
%які мають більше чотирьох різних натуральних дільників

\par
%end
\vspace{0.9cm}
%begin
148) \space 
З послідовності цілих вилучили всі числа, у сімковому запису яких нема непарних цифр. Знайти та вивести першу найдовшу підпослідовність залишку, що є власною передгорою та має довжину не менше двох.
%власною передгорою
%у сімковому запису яких нема непарних цифр

\par
%end
\vspace{0.9cm}
%begin
149) \space 
З послідовності цілих вилучили всі числа, у сімковому запису яких нема цифри 2. Знайти та вивести першу найдовшу підпослідовність залишку, що є зростаючою серією та має довжину не менше двох.
%зростаючою серією
%у сімковому запису яких нема цифри 2

\par
%end
\vspace{0.9cm}
%begin
150) \space 
З послідовності цілих вилучили всі числа, у трійковому запису яких нема парних цифр. Знайти та вивести першу найдовшу підпослідовність залишку, що є власною післягорою та має довжину не менше двох.
%власною післягорою
%у трійковому запису яких нема парних цифр

\par
%end
\vspace{0.9cm}
\end{document}
